 %%%%%%%%%%%%%%%%%%%%%%%%%%%%%%%%%%%%%%%%%
% Customized the following:
% Medium Length Professional CV
% LaTeX Template
% Version 3.0 (December 17, 2022)
% author:
% Trey Hunner (http://www.treyhunner.com/)
% License:
% CC BY-NC-SA 4.0 (https://creativecommons.org/licenses/by-nc-sa/4.0/)
%%%%%%%%%%%%%%%%%%%%%%%%%%%%%%%%%%%%%%%%%

%----------------------------------------------------------------------------------------
%	PACKAGES AND OTHER DOCUMENT CONFIGURATIONS
%----------------------------------------------------------------------------------------

\documentclass[
	a4paper, % Uncomment for A4 paper size (default is US letter)
	12pt, % Default font size, can use 10pt, 11pt or 12pt
]{resume} % Use the resume class


\title{Resume.PooriaAskariM}

%------------------------------------------------

\name{Pooria Askari Moqaddam} % Your name to appear at the top
\subtitle{Software Engineer}  % Your subtitle description to appear at the top


\email{\href{mailto:pooriaaskarim@gmail.com}{pooriaaskarim@gmail.com}}
\phone{\href{tel:+989337330917}{+98 933 733 0917}}
\address{Shiraz, Iran}
\linkedin{\href{https://www.linkedin.com/in/pooriaaskarim/}{linkedin.com/in/pooriaaskarim}}
\github{\href{https://github.com/pooriaaskarim}{github.com/pooriaaskarim}}

%----------------------------------------------------------------------------------------
\color{darkgray}
\begin{document}
\makeheader
\raggedright

%----------------------------------------------------------------------------------------
%	Profile
%----------------------------------------------------------------------------------------

\begin{rSection}{Profile}
\smallskip
I am a cross-platform developer with a strong inclination toward systems thinking---I find genuine satisfaction in understanding how things work at a deeper level and finding ways to make them work better. I have led small development teams and learned firsthand that great software is rarely a solo endeavor; the quality of what a team builds is inseparable from the quality of how it communicates. I care about both sides of that equation: writing clean, well-reasoned code and fostering the kind of engineering culture where good ideas can surface and thrive. I still have a lot to learn, and that is precisely what keeps me motivated.
\end{rSection}

%----------------------------------------------------------------------------------------
%	Brief Overview
%----------------------------------------------------------------------------------------

\begin{rSection}{Brief Overview}
\smallskip
My \hyperlink{experience}{\textbf{professional work}} spans roughly five years in software---real-time financial data systems, CRM applications, large-scale ERP platforms, and Linux infrastructure. Before that, I spent several years as a hardware technician and network administrator: component-level diagnostics and repairs, local network deployment, and server provisioning. That background still informs how I think about software.

In parallel, I have been building \hyperlink{opensource}{\textbf{open source developer tooling}} — both projects accumulated over years of real engineering work before taking their current shape:
\begin{list}{$\pdot$}{\leftmargin=12pt \itemsep=0.1em \parsep=0pt \topsep=2pt}
    \item \textbf{Logd} --- a hierarchical logging engine for Dart/Flutter. Loggers are organized in a dot-separated tree with lazy config resolution, version-based cache invalidation, and atomic multi-line buffering. A fail-safe internal pipeline ensures logging failures are never surfaced as application errors.
    \item \textbf{Flutter Notification Queue} --- a full overlay-based notification system for Flutter. Manages notifications as stateful, interactive entities through a channel-based queue architecture with 8 positional slots, drag-to-relocate, backpressure handling, and smooth gesture interactions.
\end{list}

A full breakdown of my \hyperlink{skills}{\textbf{technical skills}} follows directly below, with \hyperlink{education}{\textbf{academic background}} at the end---the latter spanning electrical engineering and computer science, a combination that largely explains the rest of this document.
\end{rSection}

%----------------------------------------------------------------------------------------
%	Skills SECTION
%----------------------------------------------------------------------------------------

\begin{rSection}{Skills}
    \hypertarget{skills}{}
    \smallskip
    \renewcommand{\arraystretch}{1.5}%
	\begin{tabular}{@{} >{\bfseries}l @{\hspace{8pt}} l @{}}
            Technologies &
                \bigChip{Dart}\bigChip{Python}\bigChip{Kotlin}\bigChip{Docker}\bigChip{Linux} \\
            
            Tools \& APIs &
                \bigChip{Flutter}\bigChip{BLoC}\bigChip{Riverpod}\bigChip{GetX}\bigChip{RxDart} \\
                &  \bigChip{Provider}\bigChip{Firebase APIs}\bigChip{OpenStreetMaps API} \\   
                &  \bigChip{Google Console APIs}\bigChip{Azure}\bigChip{Figma} \\

            Practices &
                \bigChip{Domain-Driven Design (DDD)}\bigChip{Software Architecture} \\
                &  \bigChip{Design Patterns}\bigChip{Clean Code}\bigChip{SOLID}\bigChip{OOP} \\
                &  \bigChip{Agile/Scrum}\bigChip{CI/CD} \\
            Databases &
                \bigChip{Redis}\bigChip{Hive}\bigChip{Isar}\bigChip{MySQL} \\
	\end{tabular}
\end{rSection}

\newpage

%----------------------------------------------------------------------------------------
%	WORK EXPERIENCE SECTION
%----------------------------------------------------------------------------------------

\begin{rSection}{Professional Experience}
\hypertarget{experience}{}
\smallskip
%------------------------------------------------
    \begin{rSubsection}{Rahandaaz}{2024 - 2025}{Software Developer}{Shiraz, Iran}
		\item Led a cross-functional team of developers, designers, and business stakeholders, translating complex requirements into modular domain architectures using \textbf{Domain-Driven Design (DDD)} principles.
		\item Redesigned codebase architecture and developer workflows to resolve cross-team communication bottlenecks, aligning product UX designs with engineering practices.
        \\ \chip{Flutter} \chip{GetX} \chip{Hive} \chip{Kotlin} \chip{DDD} \chip{Figma} \chip{Firebase} \chip{Agile/Scrum}
	\end{rSubsection}

	\begin{rSubsection}{Didehban Tejarat}{2023 - 2024}{Front-end Developer}{Shiraz, Iran}
		\item Led front-end development of a \textbf{CRM application}, collaborating with cross-functional product teams to deliver \textbf{responsive interfaces} and localized components.
		\item Developed and maintained inter-team \textbf{code quality guidelines} to streamline onboarding and code reviews.
		\item Automated containerized staging environments and deployment pipelines using \textbf{Docker} and \textbf{GitHub Actions} workflows, streamlining release cycles and eliminating manual configuration.
        \\ \chip{Flutter} \chip{GetX} \chip{Hive} \chip{OpenStreetMaps} \chip{Firebase} \chip{Docker}
	\end{rSubsection}

	\begin{rSubsection}{Chartix}{2022 - 2023}{Software Developer}{Shiraz, Iran}
		\item Developed the client-side ecosystem, engineering a native \textbf{Android application}, a responsive \textbf{PWA}, and an interactive \textbf{Telegram Bot} to unify access to financial services.
		\item Engineered a reliable data collection and caching pipeline utilizing \textbf{Redis}, implementing service health monitoring to ensure high availability and prevent data loss during high-volume market activity.
		\item Optimized UI state-management and rendering layers in Flutter to achieve \textbf{high-FPS chart interactions} (zooming, panning) and eliminate UI thread blocking from high-frequency WebSocket ticks.
        \\ \chip{Flutter} \chip{Python} \chip{BloC} \chip{Hive} \chip{Socket} \chip{Redis} \chip{TelegramAPI}
	\end{rSubsection}

	\begin{rSubsection}{FepCo}{2021 - 2022}{Junior Front-end Developer}{Shiraz, Iran}
		\item Practiced \textbf{Agile/Scrum} methodologies in a collaborative environment, developing a modular, large-scale \textbf{ERP software}.
		\item Developed proficiency in \textbf{Dart} and \textbf{Flutter} architectures, implementing \textbf{Clean Code}, \textbf{Domain-Driven Design (DDD)} principles, \textbf{state/route management}, and \textbf{localization}.
		\item Collaborated with design teams on \textbf{UI/UX workflows} to optimize user experience and layout fidelity.
        \\ \chip{Flutter} \chip{Dart} \chip{GetX} \chip{Clean Code} \chip{DDD} \chip{Scrum} \chip{UI/UX}
	\end{rSubsection}


	\begin{rSubsection}{Diakoo}{2019 - 2020}{System Administrator}{Shiraz, Iran}
		\item Configured and maintained \textbf{Linux/Windows server} environments optimized for hosting \textbf{Java-based} enterprise applications.
		\item Provisioned and maintained \textbf{VoIP} telecommunication servers and directory services.
		\item Collaborated with developers to design and implement a secure \textbf{video streaming service} with token-based access control and download protection.
        \\ \chip{Linux} \chip{Windows Server} \chip{VoIP} \chip{Bash} \chip{Systems Administration} \chip{Java} \chip{Video Streaming}
	\end{rSubsection}

	\begin{rSubsection}{IT Support \& Network Services}{2014 - 2019}{Hardware Technician \& Network Administrator}{Shiraz, Iran}
		\item Performed \textbf{component-level diagnostics}, motherboard repairs, and hardware upgrades for desktops, laptops, network equipment, and power supply systems.
		\item Deployed local \textbf{network infrastructures}, configured active switches/routers, and maintained power backups (\textbf{UPS}).
        \\ \chip{Hardware Diagnostics} \chip{Component Repair} \chip{Network Configuration} \chip{Troubleshooting}
	\end{rSubsection}
%------------------------------------------------
\end{rSection}



%----------------------------------------------------------------------------------------
%	OPEN SOURCE PROJECTS SECTION
%----------------------------------------------------------------------------------------

\begin{rSection}{Open Source Projects}
\hypertarget{opensource}{}
\smallskip
%------------------------------------------------
    \begin{rSubsection}{Logd}{2025 - Present}{Creator}{BSD 3-Clause 
		\\ {\footnotesize\textcolor{bluegray}{\faGithubAlt~\href{https://github.com/pooriaaskarim/logd}{pooriaaskarim/logd} $\cdot$ \faCube~\href{https://pub.dev/packages/logd}{pub.dev/packages/logd}}}}
		\item[] {\small\textit{A Dart/Flutter logging library built on the conviction that log records are structured data, deserving a proper IR, a real layout engine, and zero pressure on the garbage collector.}}
		\item Modeled the logging pipeline as a two-stage system: a semantic IR (\textbf{LogDocument}) that captures structure, followed by a physical layout pass (\textbf{TerminalLayout}) that resolves cell-aware multi-line alignment before any output is written.
		\item Designed a \textbf{LIFO object pool} (Arena) covering all IR node types and native memory buffers---eliminating per-log allocations and GC churn under high-frequency streams.
		\item Implemented a \textbf{hybrid execution mode}: Object Mode preserves the full decorator pipeline via tree traversal; Streaming Mode writes directly to native Binary IR buffers for a measured 20$\times$ throughput gain.
		\item Integrated \textbf{Dart FFI} to dispatch serialized Binary IR packets to a background isolate, with backpressure capping at 200 in-flight packets and a 128\,MB native memory ceiling.
        \\ \chip{Semantic IR} \chip{Arena Object Pool} \chip{Binary IR / FFI} \chip{Hybrid Rendering} \chip{Terminal Layout Engine}
	\end{rSubsection}

    \begin{rSubsection}{Flutter Notification Queue}{2025 - Present}{Creator}{BSD 3-Clause 
		\\{\footnotesize\textcolor{bluegray}{\faGithubAlt~\href{https://github.com/pooriaaskarim/flutter_notification_queue}{pooriaaskarim/flutter\_notification\_queue} $\cdot$ \faCube~\href{https://pub.dev/packages/flutter_notification_queue}{pub.dev/packages/flutter\_notification\_queue}}}}
		\item[] {\small\textit{A Flutter overlay system built on the premise that notifications are live agents---capable of competing for priority, responding to gestures, and yielding under backpressure---not ephemeral banners.}}
		\item Architected a \textbf{QueueCoordinator} as the lifecycle bridge between logical queue operations and Flutter's \texttt{OverlayPortal}, using a startup-mailbox pattern to handle pre-mount enqueue requests without race conditions.
		\item Implemented a \textbf{modular gesture FSM} where each interaction (drag-to-dismiss, drag-to-relocate, reorder) is an independent plugin, routed by state machine rather than nested conditional logic.
		\item Exposed a \textbf{broadcast event stream} (queued, dismissed, relocated, reordered, overflowed) that decouples observers from library internals and supports multiple simultaneous listeners.
		\item Built queue overflow as a \textbf{first-class event} rather than a silent drop---backpressure propagates to the caller, keeping the notification lifecycle fully observable.
        \\ \chip{FSM Gesture Engine} \chip{OverlayPortal} \chip{Reactive Event Bus} \chip{Backpressure} \chip{Deterministic Testing}
	\end{rSubsection}

    % \begin{rSubsection}{SarvMD}{2026 - Present}{Creator}{BUSL-1.1 \enspace {\footnotesize\textcolor{bluegray}{\faGithubAlt~\href{https://github.com/pooriaaskarim/sarvmd}{pooriaaskarim/sarvmd}}}}
	% 	\item Architected a zero-dependency \textbf{vector engraving engine} in Dart for a blank-canvas sheet music editor.
	% 	\item Implemented dynamic \textbf{page-layout calculations} and system indentation rules conforming to standard sheet music engraving conventions.
	% 	\item Designed an interactive editor UI using \textbf{custom-painted canvas} components and reactive state management.
	% 	\item Built vector export modules (\textbf{SVG/PDF}) and visual alignment guides to verify document layouts.
    %     \\ \chip{Vector Layout Engine} \chip{Gouldian Engraving} \chip{SVG/PDF Exporters} \chip{State Management} \chip{Custom Canvas}
	% \end{rSubsection}
%------------------------------------------------
\end{rSection}



%----------------------------------------------------------------------------------------
%	EDUCATION SECTION
%----------------------------------------------------------------------------------------
\begin{rSection}{Education}
	\hypertarget{education}{}
	\smallskip
	\textbf{Islamic Azad University} \hfill \textit{2018} \\ 
	Bachelor's Degree in Computer Software Technology \& Engineering \\
	\smallskip
    \textbf{Arak University of Technology} \hfill \textit{2015} \\ 
	Undergraduate Degree in Electrical Power Transmission / Installation
\end{rSection}

\end{document}
